% Generated draft: requires mathematical and source review.
\documentclass[11pt]{article}
\usepackage[margin=25mm]{geometry}
\usepackage{fontspec}
\setmainfont{Latin Modern Roman}
\usepackage{amsmath,amssymb,mathtools,unicode-math}
\setmathfont{Latin Modern Math}
\usepackage{enumitem,needspace,xurl,hyperref}
\setlength{\parindent}{0pt}
\setlength{\parskip}{5pt}
\setlength{\emergencystretch}{3em}
\newcommand{\N}{\mathbb{N}}
\newcommand{\Z}{\mathbb{Z}}
\newcommand{\Q}{\mathbb{Q}}
\newcommand{\R}{\mathbb{R}}
\newcommand{\C}{\mathbb{C}}
\newcommand{\F}{\mathbb{F}}
\newcommand{\A}{\mathbb{A}}
\newcommand{\PP}{\mathbb{P}}
\newcommand{\E}{\mathbb{E}}
\newcommand{\eps}{\varepsilon}
\newcommand{\dd}{\,\mathrm{d}}
\begin{document}
\section*{UnsolvedMath ID: 11000079}
\textbf{Source problem code:} catalogue entry
\subsection{Definitions and mathematical statement}
Let \(\lambda(f)>1\) be the stretch factor of a pseudo-Anosov mapping class \(f\in\operatorname{Mod}_g\), characterized by scaling of its invariant measured foliations. Write
\[D_g=\inf\{\lambda(f):f\in\operatorname{Mod}_g\text{ is pseudo-Anosov}\}.\]
The problem asks to compute the source's \(L(\operatorname{Mod}_g)\) explicitly for small \(g\geq2\). Its discussion of algebraic degree and minimal dilatation suggests \(L(\operatorname{Mod}_g)=D_g\), but the defining preceding notation is absent. If \(L\) instead denotes logarithmic dilatation/Teichm\"uller translation length, the target would be \(\log D_g\); these cannot silently be identified. No restriction to orientable invariant foliations or a special mapping-class subgroup is stated.

An explicit exact computation should identify the least value in the intended convention, exhibit a pseudo-Anosov attaining it, and justify exclusion of every smaller value. A numerical approximation or upper-bound example alone does not specify the minimum.

The source describes an in-principle procedure for a fixed \(g\): bound the algebraic degrees of possible dilatations, obtain bounds/candidates from algebraic integers in that degree range, enumerate train-track descriptions and pseudo-Anosov candidates, bound when dilatations become large, and match the two lists to certify a minimum. It calls practical execution difficult even for small genus. Its notation \(\lambda(\alpha)\) is described as the ``largest root'' of an algebraic integer's minimal polynomial; this needs a convention for complex roots (largest positive real root versus maximal modulus) and admissible Perron-type candidates. The requisite finite enumeration, stopping bounds and normalization are not supplied, so this outline is retained as a proposed method rather than a validated algorithm.

The research summary reports some small-genus progress, especially genus two and three, but mixes golden/silver-ratio descriptions, a tentative \((3+\sqrt5)/2\) value, and qualified resolution claims without usable references. Those conflicting numerical claims are not inserted as the value of \(D_g\) or \(L\).
\subsection{Short English statement}
Determine exact least pseudo-Anosov dilatations in small genera, with a realizing example and a matching lower bound. The source's \(L\) normalization and the inconsistent small-genus claims in its research summary require verification.
\subsection{Sources}
\begin{itemize}
\item Local UnsolvedMath registry, permanent ID \(11000079\), exact source code \texttt{catalogue entry}.
\item Farb, \emph{Problems on Mapping Class Groups and Related Topics}, Problem 7.3, PDF page 57: \url{https://www.math.uchicago.edu/~farb/papers/mcgbook.pdf}. Registry access 2026-07-29; pdf-font-structure extraction.
\item The research summary mentions Zhang and Cho and computer-search work without precise citations; its proposed values and attribution are unverified and internally qualified.
\end{itemize}
\end{document}
